% META: {"id":"TCOMPL-ME-CO-002","chapitre":"TCOMPL-MODELES-EVOLUTION","type_objet":"corrige","exercice_ref":"TCOMPL-ME-EX-002","capacites_codes":["C4"],"sources_inspiration":[],"status":"approved","origin":"generated","status_history":["generated","audited","accepted","approved"],"release_acceptance":"RELEASE_OWNER_BATCH_ACCEPTANCE","acceptance_closure_digest":"sha256:9b3ccf9a81c5520fbb7e03b7b2d3e7bf2057a02834b6d908bf3be5f49f3b553f","acceptance_receipt_digest":"sha256:4fad86f0282e0fa4e0419cca1ad6379ae7af5a280c04a4a90880f90a1c044242"}
% BEGIN-VERIFY
% from sympy import Rational
% a, b, u0 = Rational(3, 4), 5, 2
% c = b / (1 - a)
% assert c == 20
% u = u0
% for _ in range(4):
%     u = a * u + b
% formule = c + (u0 - c) * a**4
% assert u == formule
% END-VERIFY
\begin{corrige}{TCOMPL-ME-EX-002}
\textbf{1.} $c = 0{,}75c+5 \iff 0{,}25c=5 \iff c=20$.

\textbf{2.} $v_{n+1} = u_{n+1}-20 = (0{,}75u_n+5)-20 = 0{,}75u_n - 15 = 0{,}75(u_n-20) = 0{,}75\,v_n$. La suite $(v_n)$ est géométrique de raison $0{,}75$.

\textbf{3.} $v_0 = u_0-20 = 2-20=-18$, donc $v_n = -18 \times 0{,}75^n$, soit :
\[
  u_n = 20 - 18 \times 0{,}75^n.
\]
\end{corrige}
